% ===============================================================
%  جلسه بیست‌وپنجم
% ===============================================================
\chapter{آزمون فرض: مفاهیم پایه و آزمون $z$}

در این جلسه وارد آخرین مبحث بزرگ درس می‌شویم: \textbf{آزمون فرض \lr{(Hypothesis Testing)}}. تا اینجا با داده‌ها \textbf{عدد} تخمین می‌زدیم (تخمین نقطه‌ای) یا \textbf{بازه} می‌ساختیم (تخمین بازه‌ای). اما بسیاری از پرسش‌های علمی و عملی از جنس «بله یا خیر» هستند: «آیا این دارو اثر دارد؟»، «آیا این دو متغیر به هم مربوط‌اند؟». آزمون فرض چارچوبی است برای پاسخ دادن به چنین پرسش‌هایی از روی داده‌های محدود، به شکلی \textbf{احتمالاتی} و قابل اندازه‌گیری. در این فصل مفاهیم پایه (فرضیه باطل و جایگزین، آماره آزمون، ناحیه رد، خطاهای نوع اول و دوم، مقدار \lr{p}) را با دقت می‌سازیم و نخستین آزمون، یعنی \textbf{آزمون $z$} برای میانگین با واریانس معلوم، را به طور کامل بررسی می‌کنیم.

(در فایل اسلاید این جلسه، عنوان «جلسه بیست‌وششم» درج شده است؛ در این جزوه شماره جلسه‌ها به ترتیب فایل‌های اسلاید است.)

\section{پرسش‌هایی که آزمون فرض به آن‌ها پاسخ می‌دهد}
فرض کنید قرار است به پرسش‌های زیر پاسخ دهید.

\begin{question}[همبستگی درآمد و اینترنت]
میزان درآمد سرپرست یک خانوار تصادفی را با متغیر تصادفی $X$، و ضریب نفوذ اینترنت در آن خانوار را با متغیر تصادفی $Y$ نشان می‌دهیم. لذا $X,Y$ یک توزیع مشترک خواهند داشت. می‌خواهیم بدانیم که آیا میان این دو متغیر تصادفی همبستگی \lr{(correlation)} وجود دارد یا خیر؟
\begin{answer}
به عبارت دیگر، باید میان دو گزاره زیر یکی را انتخاب کنیم:
\[
\rho_{X,Y}=0\qquad\text{یا}\qquad\rho_{X,Y}\neq0 .
\]
(در اسلاید گزاره دوم به شکل برنامه‌نویسی $\rho_{X,Y}\,!\!=0$ نوشته شده که همان «نامساوی صفر» است.)
\end{answer}
\end{question}

\begin{question}[اثر واکسن]
احتمال ابتلا به بیماری کرونا برای یک شخص متوسط پیش از دریافت هیچ نوعی از واکسن را با $P_{\text{pre}}$، و احتمال ابتلای یک شخص متوسط پس از دریافت یک نوع خاص از واکسن را با $P_{\text{post}}$ نشان می‌دهیم. می‌خواهیم بدانیم که آیا تزریق واکسن احتمال را کاهش داده یا کاملاً بی‌اثر بوده است.
\begin{answer}
قصد داریم میان دو گزاره (فرض) زیر یکی را انتخاب کنیم:
\[
P_{\text{pre}}=P_{\text{post}}\qquad\text{یا}\qquad P_{\text{post}}<P_{\text{pre}} .
\]
گزاره دوم یعنی «احتمال ابتلا پس از واکسن \textbf{کمتر} شده است». (در اسلاید این گزاره به صورت $P_{\text{pre}}<P_{\text{post}}$ نوشته شده که جهت آن با متن پرسش، یعنی «کاهش احتمال»، سازگار نیست؛ صورت درست همان است که اینجا آمده.)
\end{answer}
\end{question}

\begin{question}[اثر یک برنامه تلویزیونی]
یک مقام مسئول مدعی شده که پخش یک برنامه آموزشی خاص در صداوسیما، آمار جرم و بزهکاری در جامعه را به طرز معنی‌داری کاهش داده است. فرض کنید که میزان متوسط روزانه انجام بزهکاری پیش از پخش برنامه مذکور مقدار $\mu_0$ بوده باشد. همچنین، میزان بزهکاری روزانه بعد از پخش برنامه را با متغیر تصادفی $X$ نشان می‌دهیم. آیا میانگین $X$ از مقدار پیشین $\mu_0$ کمتر شده است یا خیر؟
\begin{answer}
\[
\E(X)=\mu_0\qquad\text{یا}\qquad\E(X)<\mu_0 .
\]
\end{answer}
\end{question}

\begin{question}[کدام توزیع؟]
درآمد یک شهروند تصادفی از کشور با متغیر تصادفی $X$ تومان در سال مدل‌سازی گردیده است. می‌دانیم که توزیع $X$ منطبق با یکی از توزیع‌های دانسته‌شده $P$ یا $Q$ است. می‌خواهیم بدانیم کدام توزیع مربوط به متغیر تصادفی $X$ است.
\begin{answer}
\[
X\sim P\qquad\text{یا}\qquad X\sim Q .
\]
\end{answer}
\end{question}

و بسیاری دیگر پرسش‌هایی از این دست\ldots{} وجه مشترک همه: \textbf{دو گزاره رقیب} درباره توزیع (یا پارامترهای توزیع) یک یا چند متغیر تصادفی، و \textbf{داده‌هایی} که باید بر اساس آن‌ها یکی را انتخاب کنیم.

\subsection{چرا چنین سؤالاتی زیاد پرسیده می‌شوند؟}
\begin{itemize}
\item \textbf{نکته مهم:} این قبیل پرسش‌ها و نحوه پاسخ‌گویی به آنان یک مفهوم پایه‌ای و بنیادین جدید در علم آمار و احتمال نیست! بلکه به استفاده از روش‌های سنتی آمار و احتمال (که تاکنون دیده‌ایم: میانگین نمونه‌ای، قضیه حد مرکزی، توزیع $t$، \ldots) در پاسخ دادن به اشکال خاص و منحصر به فردی از سؤالات اشاره دارند.
\item دلیل ظاهر شدن این اشکال خاص از سؤالات، رویکردهای غالب پژوهشی در بررسی بسیاری از مسائل و پرسش‌های مطرح در علوم طبیعی و انسانی است:
\begin{itemize}
\item برای مثال، پیش از بررسی جامع یک شکل خاص از ارتباط میان چندین متغیر، می‌خواهیم بدانیم که آیا اصلاً با یکدیگر ارتباطی دارند یا خیر؟
\item و یا در بسیاری از موارد، وجود یا عدم وجود یک پدیده یا شکل خاصی از رابطه، از چرایی و مکانیزم عملکردی آن مهم‌تر است. (مثلاً برای تأیید یک دارو، ابتدا باید ثابت شود که \textbf{اثر دارد}؛ اینکه دقیقاً چگونه اثر می‌کند پرسش بعدی است.)
\end{itemize}
\end{itemize}

\subsection{چالش اصلی}
\begin{itemize}
\item در تمامی پرسش‌های مطرح‌شده، سؤالاتی در مورد توزیع یا برخی آماره‌های یک متغیر تصادفی، و یا آماره‌های مشترک دو یا چند متغیر تصادفی پرسیده شده است.
\item در صورتی که توزیع متغیر(های) تصادفی مذکور معلوم و دانسته شده باشد، پاسخ به پرسش‌های پیشین سرراست بوده و نیازمند برخی محاسبات ساده است. منظور، محاسبه \lr{Pearson correlation} (منظور رابطه $\rho_{X,Y}=\frac{\Cov(X,Y)}{\sigma_X\sigma_Y}$) میان دو متغیر تصادفی $X,Y$ و امثالهم است.
\item \textbf{چالش اصلی اینجاست که توزیع واقعی متغیر(های) تصادفی مورد بحث معلوم و دانسته شده نیست، و صرفاً تعدادی نمونه از توزیع‌ها در اختیار ما قرار گرفته است.}
\end{itemize}
مثلاً حتی اگر واقعاً $\rho_{X,Y}=0$ باشد، همبستگی نمونه‌ای ۱۰۰ خانوار تقریباً هیچ‌گاه دقیقاً صفر نمی‌شود (مثلاً $0.07$ می‌شود). آیا $0.07$ «به اندازه کافی» از صفر دور است که بگوییم همبستگی وجود دارد، یا این مقدار صرفاً حاصل تصادف نمونه‌گیری است؟ این همان سؤالی است که آزمون فرض پاسخ می‌دهد.

\subsection{راه‌حل}
\begin{itemize}
\item لذا نمی‌توان بدین گونه سؤالات جواب \textbf{قطعی} و \textbf{یقینی} داد؛ چون آماره‌ها و یا ویژگی‌های خواسته‌شده از توزیع‌ها (در اکثریت قریب به اتفاق موارد) از روی تعداد محدودی نمونه قابل حصول نیستند.
\item راه‌حل در این گونه موارد، ارائه یک جواب \textbf{احتمالاتی} است. یعنی مثلاً بگوییم فرض مذکور با (حداقل) چه احتمالی برقرار بوده و یا با (حداکثر) چه احتمالی رد می‌شود.
\end{itemize}

\begin{note}
\notetitle \textbf{تمثیل دادگاه.} منطق آزمون فرض شبیه دادگاه است. متهم «بی‌گناه فرض می‌شود مگر خلافش ثابت شود» (این همان \textbf{فرضیه باطل} است). دادستان مدارک (\textbf{داده‌ها}) ارائه می‌کند. اگر مدارک با فرض بی‌گناهی «بسیار بعید» باشند (فراتر از شک معقول)، فرض بی‌گناهی رد و حکم محکومیت صادر می‌شود. اما اگر مدارک کافی نباشند، حکم «تبرئه» است، که به معنی \textbf{اثبات بی‌گناهی نیست}، بلکه فقط یعنی «شواهد کافی برای محکومیت وجود نداشت». دو نوع اشتباه هم ممکن است: محکوم کردن بی‌گناه (خطای نوع اول) و تبرئه کردن گناهکار (خطای نوع دوم). همه این مفاهیم را در ادامه دقیق می‌کنیم.
\end{note}

\section{رهیافت ریاضیاتی و مقدمات لازم}

\subsection{مراحل انجام یک آزمون فرض}
\begin{enumerate}
\item \textbf{ایجاد یک فرض \lr{(Hypothesis)}} در مورد شکل، مقدار یک پارامتر و یا \ldots{} برای یک توزیع نامعلوم.
\item \textbf{جمع‌آوری تعدادی نمونه} (ترجیحاً \lr{i.i.d.}) از توزیع مذکور.
\item \textbf{محاسبه یک آماره نمونه‌ای \lr{(Sample Statistics)}} مربوط به توزیع، از روی نمونه‌های مشاهده‌شده.
\item \textbf{مقایسه آماره نمونه‌ای محاسبه‌شده با مقدار (یا مقادیر ممکن) فرض‌شده درباره آن:} در این مرحله، می‌بایست یک احتمال رد یا قبولی برای فرض بر اساس مشاهدات آماری ارائه شود.
\end{enumerate}
ترتیب مراحل مهم است: فرض‌ها (و حتی آستانه تصمیم) باید \textbf{پیش از} دیدن داده‌ها مشخص شوند. اگر ابتدا داده‌ها را ببینیم و سپس فرضی بسازیم که با آن‌ها جور درآید، احتمالات محاسبه‌شده دیگر معنای درستی ندارند.

\subsection{تعاریف لازم برای انجام آزمون فرض}

\begin{definition}[فرضیه ساده و مرکب]
\begin{itemize}
\item \textbf{فرضیه ساده \lr{(Simple Hypothesis)}:} به فرضیه‌ای اطلاق می‌گردد که توزیع یک یا چند متغیر تصادفی مورد سؤال را به صورت \textbf{دقیق} مشخص می‌کند.
\item \textbf{فرضیه مرکب \lr{(Composite Hypothesis)}:} به فرضیه‌ای اطلاق می‌گردد که توزیع متغیر(های) تصادفی مورد بحث را به صورت دقیق مشخص نکرده، و صرفاً نواحی خاصی را برای پارامترهای آنان مشخص می‌کند.
\end{itemize}
\end{definition}
مثال: اگر $X_i\sim\mathcal N(\mu,1)$ باشد، فرضیه «$\mu=3$» ساده است (توزیع دقیقاً $\mathcal N(3,1)$ می‌شود)، اما فرضیه «$\mu\neq3$» یا «$\mu<3$» مرکب است (بی‌نهایت توزیع ممکن را در بر می‌گیرد). در پرسش چهارم بالا هر دو فرضیه $X\sim P$ و $X\sim Q$ ساده‌اند.

\begin{definition}[آماره \lr{(Statistic)}]
به مقداری اطلاق می‌گردد که \textbf{تماماً} از روی نمونه‌ها به دست آمده و عصاره اطلاعات موجود در نمونه‌ها را (درباره پرسش مورد نظر) در خود دارد.
\end{definition}
«تماماً از روی نمونه‌ها» یعنی برای محاسبه آماره نباید به پارامتر نامعلوم نیاز باشد. مثلاً $\bar X_n$ یک آماره است، اما $\bar X_n-\mu$ (با $\mu$ نامعلوم) آماره نیست. در آزمون فرض، آماره می‌تواند به مقادیر \textbf{فرض‌شده} (مثل $\mu_0$ که معلوم است) وابسته باشد.

\begin{definition}[فرضیه باطل و فرضیه جایگزین]
\begin{itemize}
\item \textbf{فرضیه باطل \lr{(Null Hypothesis)} یا $H_0$:} به فرضیه‌ای گفته می‌شود که محقق \lr{(investigator)} معمولاً قصد رد آن را دارد. یعنی در حالت عادی، تصور عامه مبنی بر درست بودن این فرض است ولی محقق قصد دارد خلاف آن را نشان دهد. فرضیه باطل معمولاً یک \textbf{فرضیه ساده} است.
\item \textbf{فرضیه جایگزین \lr{(Alternative Hypothesis)} یا $H_1$:} فرضیه مکمل فرضیه باطل را فرضیه جایگزین می‌گویند؛ یعنی هر فرضیه‌ای به غیر از فرضیه باطل. محقق قصد دارد درستی این فرضیه را در مقابل فرضیه باطل نمایان سازد. فرضیه جایگزین معمولاً یک \textbf{فرضیه مرکب} است.
\end{itemize}
\end{definition}

به عبارت دیگر:
\begin{center}
\renewcommand{\arraystretch}{1.45}
\begin{tabular}{>{\raggedright\arraybackslash}p{5.6cm}|>{\raggedright\arraybackslash}p{5.6cm}}
فرضیه باطل $H_0$ & فرضیه جایگزین $H_1$ یا $H_A$\\\hline
فرضیه‌ای در تقابل با هدف مطالعه & فرضیه‌ای که هدف اصلی مطالعه است\\
فرضیه‌ای که می‌خواهید ادعا کنید \textbf{غلط} است & فرضیه‌ای که می‌خواهید ادعا کنید \textbf{درست} است\\
معمولاً ساده: «اثری وجود ندارد»، «تفاوتی نیست» & معمولاً مرکب: «اثر وجود دارد»
\end{tabular}
\end{center}

\textbf{چرا $H_0$ معمولاً ساده است؟} چون کل محاسبات آزمون \textbf{با فرض درستی $H_0$} انجام می‌شود: باید بتوانیم بگوییم «اگر $H_0$ درست باشد، آماره آزمون چه توزیعی دارد؟». اگر $H_0$ ساده باشد، توزیع داده‌ها و در نتیجه توزیع آماره کاملاً مشخص است و این احتمالات قابل محاسبه‌اند. در پرسش‌های بالا: $\rho=0$، $P_{\text{pre}}=P_{\text{post}}$ و $\E X=\mu_0$ همگی فرضیه‌های «بی‌اثری» هستند که نقش $H_0$ را دارند.

\begin{definition}[آزمون آماری \lr{(Statistical Test)}]
به پروسه تبدیل نمونه‌های مشاهده‌شده به یک آماره، و سپس ارائه یک گزاره احتمالاتی در مورد رد یا پذیرش فرضیه باطل بر مبنای مقدار و توزیع دقیق/تقریبی آماره اطلاق می‌گردد.
\end{definition}
\begin{itemize}
\item در طی این پروسه، یک آماره مانند $z=z(X_1,X_2,\ldots,X_n)$ از روی دادگان ساخته شده و بر حسب مقدار آن یک گزاره غیریقینی در مورد فرضیه باطل $H_0$ (یا جایگزین $H_1$) گفته خواهد شد.
\item این کار در عمده موارد از طریق مقایسه مقدار آماره با یک آستانه خاص (که بعضاً به آن \lr{critical value} نیز اطلاق می‌گردد) انجام می‌پذیرد.
\item آماره‌های مورد اشاره در عمده موارد توابعی چون میانگین یا واریانس نمونه‌ای، همبستگی نمونه‌ای و امثالهم هستند.
\end{itemize}

\subsection{ناحیه رد/قبولی و مقدار \lr{p}}
\begin{definition}[ناحیه قبولی و ناحیه رد \lr{(Acceptance/Rejection Region)}]
\begin{itemize}
\item \textbf{ناحیه قبولی \lr{(Acceptance Region)}} به زیرمجموعه‌ای از مقادیر آماره آزمون اطلاق می‌گردد که به ازای آنان فرضیه باطل $H_0$ قبول می‌شود.
\item \textbf{ناحیه رد \lr{(Rejection Region)}} به زیرمجموعه مکمل ناحیه قبولی گفته می‌شود؛ یعنی مقادیری از آماره آزمون که به ازای آن فرضیه باطل $H_0$ رد شده و فرضیه جایگزین $H_1$ قبول می‌گردد.
\item مرز بین این دو ناحیه را معمولاً \textbf{مقدار بحرانی \lr{(Critical Value)}} مشخص می‌کند.
\end{itemize}
\end{definition}

\begin{warning}
\warntitle «قبول $H_0$» در اصطلاح آماری به معنی «اثبات درستی $H_0$» نیست، بلکه یعنی «داده‌ها شواهد کافی برای رد آن فراهم نکردند» (مانند تبرئه در دادگاه). به همین دلیل بسیاری از آماردانان به جای «قبول $H_0$» می‌گویند «$H_0$ رد نشد» \lr{(fail to reject)}. دلیل: اگر $H_0$ ساده باشد (مثلاً $\mu=\mu_0$)، مقادیر بسیار نزدیک (مثلاً $\mu=\mu_0+10^{-6}$) با هیچ تعداد معقولی از داده قابل تمایز نیستند؛ پس هیچ‌گاه نمی‌توان «ثابت کرد» که $\mu$ دقیقاً برابر $\mu_0$ است.
\end{warning}

\subsection{خطاهای نوع اول و دوم \lr{(Type I and II Errors)}}
نتیجه یک آزمون آماری یا رد فرضیه باطل است یا عدم رد آن. از طرف دیگر، در واقعیت (که ما نمی‌دانیم) یا $H_0$ درست است یا $H_1$. پس چهار حالت ممکن داریم:

\begin{center}
\renewcommand{\arraystretch}{1.5}
\begin{tabular}{c|c|c}
& در واقعیت $H_1$ درست است (سیگنال هست) & در واقعیت $H_0$ درست است (سیگنال نیست)\\\hline
تصمیم: رد $H_0$ (آشکارسازی) & مثبت درست \lr{(TP)} & مثبت کاذب \lr{(FP)}: \textbf{خطای نوع اول}\\
تصمیم: عدم رد $H_0$ & منفی کاذب \lr{(FN)}: \textbf{خطای نوع دوم} & منفی درست \lr{(TN)}
\end{tabular}
\end{center}
دو نسبت مهم (با شمارش تعداد هر حالت در تکرارهای زیاد آزمون):
\[
\text{TPF}=\text{حساسیت}=\text{توان}=\frac{TP}{TP+FN},\qquad
\text{FPF}=1-\text{ویژگی}=1-CL=\frac{FP}{FP+TN}.
\]
ستون اول فقط حالت‌هایی را می‌شمارد که $H_1$ درست است و ستون دوم فقط حالت‌هایی که $H_0$ درست است.

\begin{definition}[سطح معناداری، توان آزمون]
\begin{itemize}
\item احتمال خطای نوع اول را با $\alpha$ نشان می‌دهیم و \textbf{سطح معناداری \lr{(significance level)}} آزمون می‌نامیم:
\[
\alpha=\Prob\big(\text{رد}\ H_0\ \given\ H_0\ \text{درست است}\big).
\]
\item احتمال خطای نوع دوم را با $\beta$ نشان می‌دهیم و $1-\beta$ را \textbf{توان \lr{(power)}} یا \textbf{حساسیت} آزمون می‌نامیم:
\[
1-\beta=\Prob\big(\text{رد}\ H_0\ \given\ H_1\ \text{درست است}\big).
\]
\end{itemize}
\end{definition}
در جدول اسلاید، $1-\alpha$ با $CL$ (سطح اطمینان) و «ویژگی \lr{(specificity)}» یکی است: احتمال اینکه وقتی واقعاً سیگنالی نیست، آزمون هم چیزی اعلام نکند. اگر $H_1$ مرکب باشد، توان به مقدار واقعی پارامتر بستگی دارد و یک \textbf{تابع} است، نه یک عدد.

\textbf{مصالحه میان دو خطا.} با ثابت بودن تعداد داده‌ها نمی‌توان هر دو خطا را هم‌زمان دلخواه کوچک کرد. اگر ناحیه رد را کوچک کنیم (آستانه سخت‌گیرانه‌تر)، $\alpha$ کم می‌شود ولی آزمون اثرهای واقعی را هم کمتر آشکار می‌کند و $\beta$ زیاد می‌شود؛ و برعکس. (در مثال دادگاه: سخت‌گیری بیشتر در محکومیت، تعداد بی‌گناهان محکوم را کم و تعداد گناهکاران تبرئه‌شده را زیاد می‌کند.) تنها راه کم کردن هم‌زمان هر دو خطا، \textbf{افزایش داده‌ها} است. قرارداد رایج: $\alpha$ را کوچک و ثابت (مثلاً $0.05$) انتخاب می‌کنیم، چون خطای نوع اول (ادعای کشفی که وجود ندارد) را معمولاً جدی‌تر می‌دانیم، و سپس در میان آزمون‌هایی با این $\alpha$، آزمونی با بیشترین توان را می‌جوییم.

\begin{center}
\begin{tikzpicture}
\end{tikzpicture}\\
{\small توزیع آماره آزمون تحت $H_0$ (چپ) و تحت یک حالت از $H_1$ (راست)، و یک آستانه یک‌طرفه $\tau$. ناحیه قرمز (سمت راست آستانه زیر منحنی $H_0$) احتمال خطای نوع اول $\alpha$ و ناحیه نارنجی (سمت چپ آستانه زیر منحنی $H_1$) احتمال خطای نوع دوم $\beta$ است. جابه‌جا کردن آستانه یکی را کم و دیگری را زیاد می‌کند.}
\end{center}

\subsection{مقدار \lr{p} \lr{(p-value)}}
\begin{definition}[مقدار \lr{p}]
\textbf{مقدار \lr{p}} احتمال این است که با فرض درست بودن فرضیه باطل $H_0$، کماکان داده‌هایی با میزان پرت بودن (حداقل به اندازه) نمونه‌های مشاهده‌شده $X_1,X_2,\ldots,X_n$ دیده شود که منجر به پذیرش فرضیه جایگزین $H_1$ گردند. اگر $z_{\text{obs}}$ مقدار مشاهده‌شده آماره باشد و مقادیر «دورتر از $H_0$» مقادیر بزرگ‌تر $|z|$ باشند:
\[
\text{p-value}=\Prob\big(|Z|\ge|z_{\text{obs}}|\ \given\ H_0\ \text{درست است}\big).
\]
در اسلاید این مفهوم به صورت خلاصه با عبارت $\Prob\big(H_1\ \text{is accepted}\given H_0\ \text{is correct}\big)$ نوشته شده است؛ یعنی احتمال خطای نوع اول، اگر آستانه را دقیقاً روی داده مشاهده‌شده می‌گذاشتیم.
\end{definition}

\textbf{شهود.} مقدار \lr{p} می‌گوید: «اگر واقعاً هیچ اثری وجود نداشت ($H_0$ درست بود)، چقدر احتمال داشت صرفاً از روی شانس نتیجه‌ای به این شدت (یا شدیدتر) ببینیم؟»
\begin{itemize}
\item \lr{p} کوچک (مثلاً $0.001$): داده‌ها با $H_0$ بسیار ناسازگارند؛ یا اتفاقی بسیار نادر رخ داده یا $H_0$ غلط است. شواهد قوی علیه $H_0$.
\item \lr{p} بزرگ (مثلاً $0.4$): چنین داده‌هایی تحت $H_0$ کاملاً عادی‌اند؛ شاهدی علیه $H_0$ نداریم.
\end{itemize}
\textbf{رابطه با سطح معناداری:} اگر آزمون را با سطح $\alpha$ انجام دهیم، $H_0$ رد می‌شود \textbf{اگر و تنها اگر} $\text{p-value}\le\alpha$. پس مقدار \lr{p} کوچک‌ترین سطح معناداری است که با آن، داده‌های فعلی هنوز به رد $H_0$ منجر می‌شوند.

\begin{warning}
\warntitle دو برداشت نادرست بسیار رایج:
\begin{itemize}
\item مقدار \lr{p} \textbf{احتمال درست بودن $H_0$} نیست! $\text{p-value}=\Prob(\text{داده‌های این‌قدر پرت}\given H_0)$ است، نه $\Prob(H_0\given\text{داده‌ها})$. (این دو مانند $\Prob(A\given B)$ و $\Prob(B\given A)$ متفاوت‌اند؛ برای دومی به دیدگاه بیزی و احتمال پیشین $H_0$ نیاز داریم.)
\item \lr{p} کوچک لزوماً به معنای \textbf{اثر بزرگ} نیست. با $n$ بسیار بزرگ، حتی اختلافی ناچیز و بی‌اهمیت از نظر عملی هم مقدار \lr{p} بسیار کوچکی می‌دهد. «معنادار آماری» با «مهم در عمل» یکی نیست.
\end{itemize}
\end{warning}

\subsection{انواع آزمون‌های معروف}
انواع گوناگونی از آزمون‌های آماری برای کاربردهای مختلف طراحی شده‌اند. در ادامه (این جلسه و جلسه بعد) به برخی از آنان اشاره خواهد شد:
\begin{itemize}
\item آزمون‌های میانگین $z$-\lr{test} (واریانس معلوم)؛
\item آزمون‌های $t$-\lr{test} (واریانس نامعلوم)؛
\item آزمون \lr{Neyman--Pearson} (یا آزمون نسبت درست‌نمایی، \lr{Likelihood Ratio Test})؛
\item آزمون جایگشت \lr{(Permutation Test)}؛
\item آزمون بهنجارش \lr{($\chi^2$-test)}؛
\item و بسیاری دیگر\ldots
\end{itemize}

\section{آزمون میانگین نمونه‌ای با واریانس معلوم \lr{(Z-test)}}

\subsection{صورت مسئله}
فرض کنید که نمونه‌های $X_1,X_2,\ldots,X_n$ به صورت \lr{i.i.d.} از یک توزیع نامعلوم تولید شده باشند. فرض کنید که انحراف معیار توزیع دانسته شده و مقدار معلوم $\sigma$ را دارا باشد. هدف طراحی آزمون این فرضیه است که میانگین توزیع مقدار $\mu_0$ را دارا است یا خیر. مطابق زیر عمل خواهیم کرد:
\begin{itemize}
\item \textbf{فرضیه باطل $H_0$:} میانگین توزیع برابر با $\mu_0$ است، یا $\E(X)=\mu_0$.
\item \textbf{فرضیه جایگزین $H_1$:} میانگین توزیع برابر با $\mu_0$ نیست، یا $\E(X)\neq\mu_0$.
\end{itemize}
آماره زیر را به عنوان آماره آزمون انتخاب می‌کنیم:
\[
\boxed{\ z_n=\frac{\sqrt n\,\big(\bar X_n-\mu_0\big)}{\sigma}\ }
\]
شهود: $\bar X_n-\mu_0$ فاصله میانگین نمونه‌ای از مقدار فرض‌شده است. اما «بزرگ یا کوچک» بودن این فاصله نسبی است و باید با میزان نوسان طبیعی $\bar X_n$ مقایسه شود، که همان انحراف معیار $\sigma/\sqrt n$ است. پس $z_n$ می‌گوید: «میانگین نمونه‌ای چند برابر انحراف معیارش از $\mu_0$ فاصله دارد؟» (این آماره است، چون $\mu_0$ و $\sigma$ معلوم‌اند.)

\subsection{رفتار آماره تحت $H_0$ و $H_1$}
در صورت درست بودن فرض $H_0$، $z_n$ حول میانگین صفر تمرکز خواهد کرد (با واریانس واحد) و در صورت درست نبودن $H_0$، با افزایش $n\to\infty$ مقدار این آماره به صورت مجانبی با نرخ $n^{1/2}$ افزایش (یا کاهش) خواهد یافت. برای دیدن این، فرض کنید میانگین واقعی $\mu$ باشد و $z_n$ را تجزیه کنیم:
\[
z_n=\frac{\sqrt n\,(\bar X_n-\mu)}{\sigma}+\frac{\sqrt n\,(\mu-\mu_0)}{\sigma}.
\]
جمله اول میانگین صفر و واریانس $\frac{n}{\sigma^2}\cdot\frac{\sigma^2}{n}=1$ دارد. جمله دوم یک عدد ثابت (غیرتصادفی) است:
\begin{itemize}
\item اگر $H_0$ درست باشد ($\mu=\mu_0$)، جمله دوم صفر است و $\E z_n=0$، $\Var z_n=1$.
\item اگر $H_1$ درست باشد ($\mu\neq\mu_0$)، جمله دوم برابر $\frac{\sqrt n(\mu-\mu_0)}{\sigma}$ است که قدرمطلق آن مانند $\sqrt n$ رشد می‌کند؛ پس $z_n$ به سمت $+\infty$ (اگر $\mu>\mu_0$) یا $-\infty$ (اگر $\mu<\mu_0$) می‌رود.
\end{itemize}

\textbf{توزیع دقیق یا تقریبی:}
\begin{itemize}
\item اگر $X_i$ها که نمونه‌های محقق‌شده $X$ هستند توزیع گاوسی داشته باشند، همواره خواهیم داشت $z_n\sim\mathcal N(?,1)$، که مقدار علامت سؤال به برقرار بودن فرضیه باطل $H_0$ یا فرضیه جایگزین بستگی دارد. \textbf{چرا؟} چون $\bar X_n$ (ترکیب خطی گاوسی‌های مستقل) گاوسی است و $z_n$ تبدیلی خطی از آن است؛ با تجزیه بالا، مقدار علامت سؤال برابر $\frac{\sqrt n(\mu-\mu_0)}{\sigma}$ است، که تحت $H_0$ صفر می‌شود:
\[
z_n\sim\mathcal N\Big(\frac{\sqrt n(\mu-\mu_0)}{\sigma},\ 1\Big),\qquad\text{تحت}\ H_0:\ z_n\sim\mathcal N(0,1).
\]
\item اگر $X$ توزیع گاوسی نداشته باشد نیز، با افزایش محدود $n$، توزیع $z_n$ به یک توزیع گاوسی نزدیک و نزدیک‌تر خواهد شد (\textbf{قضیه حد مرکزی}).
\end{itemize}

\subsection{ناحیه پذیرش و مقدار بحرانی}
\begin{itemize}
\item لذا می‌توان برای رد یا قبولی فرضیه باطل $H_0$، یک ناحیه پذیرش به شکل بازه $(-\tau,\tau)$ حول صفر در نظر گرفت و اگر مقدار آماره آزمون از این ناحیه خارج بود، فرضیه جایگزین $H_1$ را قبول کرد.
\item مقدار آستانه یا \lr{critical value} مشخص‌شده با $\tau$ کاملاً سلیقه‌ای بوده و به این نکته بازمی‌گردد که محقق \lr{(investigator)} چه میزانی از شاخص اهمیت آماری \lr{(Statistical Significance Level)} را برای رد یا قبولی فرضیات خود انتخاب کرده است.
\end{itemize}
ارتباط $\tau$ و $\alpha$: احتمال خطای نوع اول برابر است با
\[
\alpha=\Prob\big(|z_n|\ge\tau\given H_0\big)=\Prob(|Z|\ge\tau)=2\big(1-\Phi(\tau)\big)
\quad\Longleftrightarrow\quad
\tau=z_{\alpha/2}=\Phi^{-1}\Big(1-\frac\alpha2\Big).
\]
پس برای $\alpha=0.05$ داریم $\tau=1.96$ و برای $\alpha=0.01$ داریم $\tau=2.576$ (همان اعداد جدول بازه اطمینان؛ این تصادفی نیست، تمرین ۳ را ببینید).

\begin{center}
\begin{tikzpicture}
\end{tikzpicture}\\
{\small آزمون $z$ دوطرفه: ناحیه عدم رد $H_0$ بازه $(-\tau,\tau)$ با احتمال $1-\alpha$ است و هر یک از دو ناحیه رد در دم‌ها احتمال $\alpha/2$ (تحت $H_0$) دارد.}
\end{center}

\begin{itemize}
\item دقت کنید که بیان قبولی/ردی فرضیه باطل توسط یک محقق چندان مهم نیست، چرا که این مورد شدیداً به ناحیه قبولی/ردی آزمون طراحی‌شده بستگی دارد که سلیقه‌ای است.
\item اما مقدار \lr{p-value} (یا دیگر معادل‌های آن مانند \lr{Significance level} یا \lr{Power}) خود را با این اعمال سلیقه شخصی بازتنظیم کرده و همواره دید درستی از قابلیت اطمینان نتیجه گزارش‌شده خواهند داد.
\end{itemize}
برای آزمون $z$ دوطرفه، مقدار \lr{p} برابر است با
\[
\text{p-value}=\Prob\big(|Z|\ge|z_{\text{obs}}|\big)=2\big(1-\Phi(|z_{\text{obs}}|)\big).
\]
مثلاً گزارش «$H_0$ در سطح ۵٪ رد نشد» اطلاعات کمی می‌دهد؛ اما گزارش «$\text{p}=0.051$» به خواننده می‌گوید که نتیجه در مرز بوده و خواننده با هر سطح معناداری دلخواه خود می‌تواند تصمیم بگیرد.

\begin{example}[آزمون $z$ دوطرفه: وزن بسته‌ها]
کارخانه‌ای ادعا می‌کند میانگین وزن بسته‌هایش $\mu_0=500$ گرم است. انحراف معیار وزن بسته‌ها از قبل معلوم و برابر $\sigma=12$ گرم است. یک بازرس $n=36$ بسته را وزن کرده و میانگین $\bar x=496.1$ گرم به دست آورده است. (الف) آزمون را در سطح $\alpha=0.05$ انجام دهید. (ب) مقدار \lr{p} را حساب کنید.
\begin{solution}
\textbf{فرضیه‌ها:} $H_0:\ \mu=500$ در برابر $H_1:\ \mu\neq500$ (بازرس نگران هر دو جهت انحراف است).

\textbf{آماره:}
\[
z_{\text{obs}}=\frac{\sqrt{36}\,(496.1-500)}{12}=\frac{6\times(-3.9)}{12}=-1.95 .
\]
(الف) $\tau=z_{0.025}=1.96$ و $|z_{\text{obs}}|=1.95<1.96$؛ پس $z_{\text{obs}}$ در ناحیه پذیرش است و $H_0$ در سطح ۵٪ \textbf{رد نمی‌شود}.

(ب)
\[
\text{p-value}=2\big(1-\Phi(1.95)\big)=2\times0.0256=0.0512 .
\]
نتیجه دقیقاً در مرز است: اگر کسی سطح $\alpha=0.06$ را انتخاب کرده بود، $H_0$ را رد می‌کرد. این همان نکته اسلاید است: «رد/قبول» به انتخاب سلیقه‌ای آستانه وابسته است، اما مقدار \lr{p} تصویر کامل‌تری می‌دهد.
\end{solution}
\end{example}

\begin{example}[توان آزمون]
در مثال قبل فرض کنید میانگین واقعی $\mu=494$ گرم باشد (۶ گرم کمتر از ادعا). احتمال اینکه آزمون (با $n=36$ و $\alpha=0.05$) این تخلف را آشکار کند چقدر است؟
\begin{solution}
طبق بخش قبل، $z_n\sim\mathcal N(d,1)$ با
\[
d=\frac{\sqrt n(\mu-\mu_0)}{\sigma}=\frac{6\times(494-500)}{12}=-3 .
\]
آزمون $H_0$ را رد می‌کند اگر $z_n<-1.96$ یا $z_n>1.96$. با نوشتن $z_n=Z+d$ که $Z\sim\mathcal N(0,1)$:
\[
\Prob(z_n<-1.96)=\Prob(Z<-1.96+3)=\Phi(1.04)=0.851,\qquad\Prob(z_n>1.96)=\Prob(Z>4.96)\approx0 .
\]
پس توان $1-\beta\approx0.851$ و احتمال خطای نوع دوم $\beta\approx0.149$ است: در حدود ۱۵٪ موارد، این آزمون تخلف ۶ گرمی را آشکار نمی‌کند. با افزایش $n$، قدرمطلق $d$ مانند $\sqrt n$ زیاد می‌شود و توان به ۱ میل می‌کند.
\end{solution}
\end{example}

\begin{summarybox}[جمع‌بندی آزمون $z$ دوطرفه]
\begin{itemize}
\item فرض‌ها: نمونه‌های \lr{i.i.d.}، $\sigma$ معلوم؛ $H_0:\ \mu=\mu_0$ در برابر $H_1:\ \mu\neq\mu_0$.
\item آماره: $z_n=\sqrt n(\bar X_n-\mu_0)/\sigma$ که تحت $H_0$ (دقیقاً برای داده گاوسی، تقریباً برای $n$ بزرگ) $\mathcal N(0,1)$ است.
\item تصمیم در سطح $\alpha$: رد $H_0$ اگر $|z_n|\ge z_{\alpha/2}$.
\item مقدار \lr{p}: $2\big(1-\Phi(|z_{\text{obs}}|)\big)$؛ رد $H_0$ اگر و تنها اگر $\text{p}\le\alpha$.
\item توان در $\mu$: $1-\beta(\mu)=\Phi(-z_{\alpha/2}-d)+1-\Phi(z_{\alpha/2}-d)$ با $d=\sqrt n(\mu-\mu_0)/\sigma$.
\end{itemize}
\end{summarybox}

\section{تمرین‌های تکمیلی}

\begin{exercise}[صورت‌بندی فرضیه‌ها]
برای هر مورد، $H_0$ و $H_1$ را بنویسید و مشخص کنید هر کدام ساده است یا مرکب. (الف) آیا یک سکه سالم است؟ (ب) آیا داروی جدید فشار خون را (به طور میانگین) پایین می‌آورد؟ (ج) آیا طول عمر یک قطعه توزیع نمایی با میانگین ۱۰۰ ساعت دارد یا نمایی با میانگین ۸۰ ساعت؟
\end{exercise}
\begin{solution}
(الف) $H_0:\ p=\frac12$ (ساده: توزیع برنولی کاملاً معلوم) در برابر $H_1:\ p\neq\frac12$ (مرکب).
(ب) با $\Delta$ = تغییر میانگین فشار خون (بعد منهای قبل): $H_0:\ \Delta=0$ (ساده، اگر بقیه پارامترها معلوم باشند) در برابر $H_1:\ \Delta<0$ (مرکب و یک‌طرفه).
(ج) $H_0:\ X\sim\mathrm{Exp}(1/100)$ در برابر $H_1:\ X\sim\mathrm{Exp}(1/80)$؛ هر دو ساده (مانند پرسش $X\sim P$ یا $X\sim Q$). برای این حالت بهترین آزمون را لم نیمن--پیرسون (جلسه بعد) می‌دهد.
\end{solution}

\begin{exercise}[آزمون $z$ یک‌طرفه: برنامه تلویزیونی]
در پرسش برنامه تلویزیونی، فرض کنید میانگین بزهکاری روزانه پیش از برنامه $\mu_0=20$ مورد و انحراف معیار روزانه $\sigma=5$ (معلوم) باشد. پس از پخش برنامه، طی $n=50$ روز میانگین $\bar x=18.6$ مشاهده شده است. با فرضیه جایگزین یک‌طرفه $H_1:\ \E X<\mu_0$، آزمون را در سطح $\alpha=0.05$ انجام دهید و مقدار \lr{p} را حساب کنید. اگر آزمون دوطرفه انجام می‌شد چه؟
\end{exercise}
\begin{solution}
$z_{\text{obs}}=\frac{\sqrt{50}\,(18.6-20)}{5}=\frac{7.071\times(-1.4)}{5}=-1.98$.

\textbf{یک‌طرفه:} چون $H_1$ فقط مقادیر کوچک‌تر را در بر دارد، فقط مقادیر بسیار \textbf{منفی} $z$ شاهدی علیه $H_0$ هستند. ناحیه رد $z\le-z_\alpha=-1.645$ است (کل $\alpha$ در یک دم). چون $-1.98<-1.645$، $H_0$ رد می‌شود. مقدار \lr{p}: $\Prob(Z\le-1.98)=\Phi(-1.98)=0.024$.

\textbf{دوطرفه:} $\text{p}=2\times0.024=0.048$ و $|z|=1.98>1.96$؛ این هم (به سختی) رد می‌کند. آزمون یک‌طرفه وقتی جهت اثر از پیش (و پیش از دیدن داده‌ها!) معلوم است، توان بیشتری دارد. (اینکه مقام مسئول کاهش بزهکاری را به این برنامه نسبت دهد، ادعای دیگری است: آزمون فقط کاهش میانگین را نشان می‌دهد، نه علت آن را.)
\end{solution}

\begin{exercise}[ارتباط آزمون دوطرفه و بازه اطمینان]
نشان دهید آزمون $z$ دوطرفه با سطح $\alpha$، فرضیه $H_0:\ \mu=\mu_0$ را رد می‌کند اگر و تنها اگر $\mu_0$ خارج از بازه اطمینان $1-\alpha$ جلسه ۲۳ باشد.
\end{exercise}
\begin{solution}
عدم رد یعنی $|z_n|<z_{\alpha/2}$:
\[
\Big|\frac{\sqrt n(\bar X_n-\mu_0)}{\sigma}\Big|<z_{\alpha/2}
\ \Longleftrightarrow\ |\bar X_n-\mu_0|<z_{\alpha/2}\frac{\sigma}{\sqrt n}
\ \Longleftrightarrow\ \bar X_n-z_{\alpha/2}\frac{\sigma}{\sqrt n}<\mu_0<\bar X_n+z_{\alpha/2}\frac{\sigma}{\sqrt n}.
\]
پس عدم رد معادل قرار داشتن $\mu_0$ در بازه اطمینان است. شهود: بازه اطمینان دقیقاً مجموعه همه مقادیر $\mu_0$ است که آزمون آن‌ها را رد نمی‌کند. در مثال وزن بسته‌ها، بازه ۹۵٪ برابر $496.1\pm1.96\times2=[492.18,\ 500.02]$ است که $500$ را (به سختی) شامل می‌شود؛ سازگار با عدم رد.
\end{solution}

\begin{exercise}[توزیع مقدار \lr{p} تحت $H_0$]
نشان دهید اگر $H_0$ درست باشد و آماره توزیع پیوسته داشته باشد، مقدار \lr{p} آزمون $z$ دوطرفه توزیع یکنواخت روی $[0,1]$ دارد. چه نتیجه‌ای می‌گیرید؟
\end{exercise}
\begin{solution}
$p=2(1-\Phi(|Z|))=g(|Z|)$ که $g$ اکیداً نزولی است. برای $u\in(0,1)$:
\[
\Prob(p\le u)=\Prob\Big(|Z|\ge\Phi^{-1}\big(1-\tfrac u2\big)\Big)=2\Big(1-\big(1-\tfrac u2\big)\Big)=u .
\]
این تابع توزیع تجمعی $\mathcal U(0,1)$ است. نتیجه: تحت $H_0$ هر مقداری از \lr{p} به یک اندازه محتمل است؛ مقدار $0.04$ همان‌قدر محتمل است که $0.74$. قانون «رد اگر $p\le\alpha$» دقیقاً احتمال خطای نوع اول $\alpha$ می‌دهد.
\end{solution}

\begin{exercise}[آزمون‌های چندگانه]
پژوهشگری ۲۰ فرضیه مستقل را هر کدام در سطح $\alpha=0.05$ آزمون می‌کند و در واقع همه فرضیه‌های باطل درست‌اند. احتمال اینکه دست‌کم یک «کشف» (رد $H_0$) گزارش کند چقدر است؟ برای اینکه این احتمال حداکثر $0.05$ شود، سطح هر آزمون را چقدر بگیرد؟
\end{exercise}
\begin{solution}
احتمال اینکه هیچ آزمونی رد نکند $0.95^{20}=0.358$ است؛ پس احتمال دست‌کم یک کشف کاذب $1-0.358=0.642$! راه‌حل ساده (\textbf{تصحیح بونفرونی}): با کران اجتماع $\Prob(\bigcup_iA_i)\le\sum_i\Prob(A_i)$ (که از اصل جمع‌پذیری و نامنفی بودن احتمال نتیجه می‌شود)، اگر هر آزمون در سطح $\alpha/m=0.05/20=0.0025$ انجام شود، احتمال دست‌کم یک خطای نوع اول حداکثر $20\times0.0025=0.05$ است (حتی بدون فرض استقلال).
\end{solution}

\begin{exercise}[تعداد نمونه برای توان دلخواه]
برای آزمون $z$ دوطرفه با $\alpha=0.05$، چند نمونه لازم است تا اگر میانگین واقعی به اندازه $\Delta=0.5\sigma$ با $\mu_0$ فاصله داشته باشد، توان آزمون (تقریباً) $0.9$ باشد؟
\end{exercise}
\begin{solution}
با $d=\sqrt n\Delta/\sigma=0.5\sqrt n>0$، توان تقریباً برابر $\Prob(Z+d>1.96)=\Phi(d-1.96)$ است (دم دیگر ناچیز است). برای توان $0.9$ باید $d-1.96=\Phi^{-1}(0.9)=1.2816$:
\[
0.5\sqrt n=3.2416\quad\Rightarrow\quad n=42.03\quad\Rightarrow\quad n=43 .
\]
فرمول کلی: $n\approx\Big(\dfrac{(z_{\alpha/2}+z_\beta)\,\sigma}{\Delta}\Big)^2$.
\end{solution}

\begin{exercise}[آیا سکه سالم است؟]
سکه‌ای ۱۰۰ بار پرتاب شده و ۶۱ بار شیر آمده است. با استفاده از قضیه حد مرکزی، $H_0:\ p=\frac12$ را در برابر $H_1:\ p\neq\frac12$ آزمون کنید.
\end{exercise}
\begin{solution}
تحت $H_0$ هر پرتاب برنولی با میانگین $\frac12$ و واریانس معلوم $\frac14$ است (اینجا $H_0$ ساده است و واریانس را هم تعیین می‌کند). پس $\sigma=0.5$ و
\[
z=\frac{\sqrt{100}\,(0.61-0.5)}{0.5}=\frac{10\times0.11}{0.5}=2.2,\qquad\text{p}=2\big(1-\Phi(2.2)\big)=0.028 .
\]
در سطح ۵٪، $H_0$ رد می‌شود (ولی در سطح ۱٪ رد نمی‌شود).
\end{solution}
